\documentclass[11pt]{article}
\usepackage[a4paper,margin=27mm]{geometry}
\usepackage{amsmath,amssymb,mathtools}
\usepackage{fontspec,unicode-math}
\setmainfont{Latin Modern Roman}
\setmathfont{Latin Modern Math}
\usepackage{xurl,hyperref}
\hypersetup{colorlinks=true,urlcolor=blue,linkcolor=blue}
\newcommand{\sref}[2]{\hyperlink{source:#1:#2}{[S#2]}}
\newcommand{\cataloguescope}[1]{\paragraph{Mathematical question.} #1}
\begin{document}
% UnsolvedMath ID 8500007; source catalogue code AMR-084-0007
\setcounter{section}{8500006}
\section{Existence of a rational Diophantine septuple}\label{8500007}
{\small\sffamily UnsolvedMath ID 8500007 · Graph Theory}
\subsection{Definitions and mathematical statement}

\paragraph{Shared notation.}$\mathbb N=\{1,2,\ldots\}$, and $\mathbb Z,\mathbb Q,\mathbb R,\mathbb C$ denote the integers, rationals, reals and complex numbers. Any other conventions are specified locally.


For $q\in\mathbb Q$, a rational $D(q)$-$m$-tuple is a set $A=\{a_1,\ldots,a_m\}\subset\mathbb Q\setminus\{0\}$ of distinct numbers such that $a_i a_j+q$ is a square in $\mathbb Q$ for every $i<j$. Elements may have either sign; zero squares are allowed.
\paragraph{Statement recorded in the supplied source.}
Does a rational $D(1)$-septuple exist, that is, can seven distinct nonzero rationals make all twenty-one expressions $a_i a_j+1$ rational squares? \sref{8500007}{2}
\subsection{Short English statement}
Can a rational Diophantine set have seven elements?
\subsection{Sources}
\begin{description}
\item[\hypertarget{source:8500007:1}{S1}] ULAM.AI and the UnsolvedMath Contributors, UnsolvedMath: A Curated Collection of Open Mathematics Problems, record 8500007 (AMR-084-0007). CC BY 4.0. \url{https://huggingface.co/datasets/ulamai/UnsolvedMath}.
\item[\hypertarget{source:8500007:2}{S2}] Andrej Dujella, Diophantine m-tuples, author-maintained survey (snapshot accessed 9 October 2026), Section 5.2, paragraph on rational sextuples and the unknown septuple. \url{https://raw.githubusercontent.com/dujella/dujella.github.io/master/ratio.html}.
\end{description}
\label{end:8500007}
\end{document}
